In \Sec{sec-realize} we assumed that every member of the set of
realised scenarios has a different sequence of events. In this section
we discuss how to relax that assumption. Our method is best explained
by an example.

Assume that we have already performed the initial partitioning
(see \Page{initial-partitioning-step}), and that
our current task is to realise the set of scenarios $S$, such
that, for every $\xi, \eta \in S$, $\Sigma_{\xi} = \Sigma_{\eta}$.
Moreover, $S = S_A \cup S_B \cup S^{\prime}$, where
\begin{itemize}
  \item
    $S_A = \{A_1, A_2, A_3\}$ and $\ESeq{A_1} = \ESeq{A_2} = \ESeq{A_3}$;
  \item
    $S_B = \{B_1, B_2\}$ and $\ESeq{B_1} = \ESeq{B_2} \neq \ESeq{A_1}$;
  \item
    every member of $S^{\prime}$ has a sequence of events that is unique in $S$.
\end{itemize}

Each projection of a DTS has a unique sequence of events, so the expression that
realises $S$ must have at least three alternatives: one for each member of
$S_A$. However, it is possible that, say, $A_2$ and $B_1$ are projections of the
same DTS. The question is how to pair the right $A$ with the right $B$.

A way to address this question is to investigate all the combinations of members
from $S_A$ and $S_B$.
More concretely, we create six sets of scenarios:
$S_1 = \{A_1, B_1\} \cup S^{\prime}$,
$S_2 = \{A_1, B_2\} \cup S^{\prime}$,
$S_3 = \{A_2, B_1\} \cup S^{\prime}$ etc.
In each of these six sets every scenario has a unique sequence of events, so we
can run our realisation algorithm separately for $S_1$, $S_2$ etc., to obtain
six scenario expressions. From each of those scenario expressions we remove
alternatives that realise only members of $S^{\prime}$, then find the best of
the remaining expressions, i.e., one with the fewest alternatives. Let us say
that it is an expression whose only alternative realises
$\{A_3, B_2, D, F\}$ (where $D,F\in S'$).
We retain this expression as a partial solution to the problem of realising
$S$, and then repeat the process for $S \setminus \{A_3, B_2, D, F\}$
to find the other alternatives of the solution for $S$.

The problem with this approach is that the number of combinations to be
investigated grows exponentially with the number of scenarios that have
identical sequences. But as long as the number of combinations is not
prohibitively large,
%a few thousand,
the approach is quite effective.
